\documentclass[11pt]{article}
\usepackage[margin=1in]{geometry}
\usepackage{amsmath,amssymb,amsthm,mathtools}
\usepackage{enumitem}
\usepackage{graphicx}
\usepackage[hidelinks]{hyperref}

\newtheorem{theorem}{Theorem}
\newtheorem{lemma}[theorem]{Lemma}
\newtheorem{proposition}[theorem]{Proposition}
\newcommand{\C}{C}
\newcommand{\Cp}{C_p}
\newcommand{\R}{\mathbb R}
\newcommand{\ev}{\operatorname{ev}}
\newcommand{\htop}{\operatorname{ht}}

\title{Positive $C_p$-homeomorphisms and positive Banach isomorphisms\\
coincide for countable compacta}
\author{Candidate full solution packet for arXiv:2601.11463, Question 4.3}
\date{21 July 2026}

\begin{document}
\maketitle

\noindent\textbf{Status.}
\texttt{candidate full, likely valid; requires human review}.

\section{Source question}
Cúth, Havelka, Rondoš, and Sarı ask the following in Question 4.3 on page 27
of~\cite{CHRS}.
\begin{center}
\includegraphics[width=\textwidth]{figures/open_problem_crop.png}
\end{center}

In transcription:
\begin{quote}
Assume that $K,L$ are countable compact spaces. Is it true that there exists a
positive linear homeomorphism of $C_p(K)$ onto $C_p(L)$ if and only if there
exists a positive isomorphism of $C(K)$ onto $C(L)$?
\end{quote}
As in~\cite{CHRS}, ``positive'' is one-sided: the displayed map is positive,
but its inverse is not required to be positive.

\section{Definitions}
For compact $K$, the vector space $\C(K)$ consists of the continuous
real-valued functions on $K$.  We write $\Cp(K)$ for this vector space with
the topology of pointwise convergence, and $\C(K)$ for the Banach space with
the supremum norm.  (The linear-topological argument is unchanged over the
complex scalars.)

For $x\in K$, let $\ev_x(f)=f(x)$.  We call a linear map
$T:\C(K)\to\C(L)$ \emph{coordinate-finite} if, for every $y\in L$, there are
finitely many $x_1,\ldots,x_r\in K$ and scalars $a_1,\ldots,a_r$ such that
\begin{equation}\label{eq:coordinate-finite}
 (Tf)(y)=\sum_{j=1}^r a_j f(x_j)\qquad(f\in\C(K)).
\end{equation}

For infinite countable compact $K$, $\htop(K)$ denotes its Cantor--Bendixson
height.  Following~\cite{CHRS}, $\Gamma(\alpha)$ is the least ordinal of the
form $\omega^\xi$ which is at least $\alpha$.

\section{Proof intuition}
The two topologies have complementary automatic-continuity mechanisms.
Uniform convergence is stronger than pointwise convergence.  Therefore, if a
linear $C_p$-homeomorphism and a candidate uniform limit are placed in a
closed-graph test, pointwise continuity identifies the limit.  The closed
graph theorem then upgrades the map and its inverse to bounded Banach-space
operators.

The reverse implication is not automatic for an arbitrary positive bounded
operator: its value at a point of $L$ could be integration against a measure
with infinite support, which is not a continuous functional for the pointwise
topology.  What resolves the problem is the particular constructive proof in
the source paper.  Its positive isomorphisms and their inverses are
coordinate-finite.  This is exactly the missing $C_p$ continuity, and the
property survives all the finite compositions and homeomorphic conjugacies in
the source classification proof.

\section{Elementary topology lemmas}
\begin{lemma}\label{lem:cp}
A linear map $T:\C(K)\to\C(L)$ is continuous from $\Cp(K)$ to $\Cp(L)$ if
and only if it is coordinate-finite.
\end{lemma}
\begin{proof}
If \eqref{eq:coordinate-finite} holds for every $y$, then every coordinate
$\ev_y\circ T$ is continuous for the pointwise topology, which is precisely
the continuity of $T$ into the product topology on $\Cp(L)$.

Conversely, fix $y\in L$.  The linear functional $\ev_y\circ T$ is continuous
on $\Cp(K)$.  Hence there are $x_1,\ldots,x_r\in K$ and $\varepsilon>0$ such
that
\[
 |f(x_j)|<\varepsilon\ (1\leq j\leq r)
 \quad\Longrightarrow\quad |(Tf)(y)|<1.
\]
If $f(x_j)=0$ for all $j$, the same implication applies to every scalar
multiple of $f$, so $(Tf)(y)=0$.  Thus $\ev_y\circ T$ factors through the
finite-dimensional coordinate map
$f\mapsto(f(x_1),\ldots,f(x_r))$.  A linear functional on its range extends
to $\R^r$, giving \eqref{eq:coordinate-finite}.
\end{proof}

\begin{lemma}\label{lem:composition}
Coordinate-finite linear maps are closed under finite composition.  A
homeomorphism $\phi:L\to K$ induces a coordinate-finite isometry
$f\mapsto f\circ\phi$, whose inverse is also coordinate-finite.
\end{lemma}
\begin{proof}
For composition, substitute each finite evaluation formula for the inner map
into the finite formula for the outer map; only a finite union of finite
supports occurs.  A composition operator has exactly one input evaluation in
each output coordinate, and the same is true for its inverse.
\end{proof}

\begin{lemma}\label{lem:closed-graph}
If $T:\Cp(K)\to\Cp(L)$ is a linear homeomorphism, then the same algebraic map
is a Banach-space isomorphism $T:\C(K)\to\C(L)$.
\end{lemma}
\begin{proof}
Suppose $f_n\to f$ uniformly on $K$ and $Tf_n\to g$ uniformly on $L$.
Uniform convergence implies pointwise convergence.  Continuity of $T$ for the
pointwise topologies gives $Tf_n\to Tf$ pointwise, while the assumed uniform
convergence gives $Tf_n\to g$ pointwise.  Hence $g=Tf$.  The norm graph of
$T$ is closed, so the closed graph theorem makes $T$ norm-bounded.  Applying
the same argument to the $C_p$-continuous inverse makes $T^{-1}$ norm-bounded.
\end{proof}

\section{The coordinate-finite content of the source construction}
We isolate the observation not stated in~\cite{CHRS}.

\begin{proposition}\label{prop:source-audit}
Let $K,L$ be infinite countable compact spaces satisfying
\begin{equation}\label{eq:height}
 \htop(K)\leq\htop(L)<\Gamma(\htop(K)).
\end{equation}
Then the constructive proof of Theorem 2.4 in~\cite{CHRS} yields a positive
isomorphism $U:\C(K)\to\C(L)$ for which both $U$ and $U^{-1}$ are
coordinate-finite.
\end{proposition}

\begin{proof}
We inspect all types of maps used in that proof.  By the
Mazurkiewicz--Sierpiński representation employed there, the spaces are first
conjugated by homeomorphisms to ordinal intervals.  Lemma~\ref{lem:composition}
handles these conjugacies.

The first nontrivial map is the operator in Proposition 2.12 of~\cite{CHRS}.
For a fixed finite $k$, its formula for $(Tf)(y)$ is either one evaluation, a
weighted sum of $k$ endpoint evaluations, or
\[
 (Tf)(\gamma)=\sum_{j=i}^k\lambda_j f(p_{i,j}(\gamma))
   +\sum_{j=1}^{i-1}\lambda_j f(\omega^\alpha j)
       \qquad(\gamma\in I_i).
\]
Thus $T$ is coordinate-finite.  The displayed inverse in that proof is also
coordinate-finite.  At the finitely many endpoints this is immediate; on the
last block it is a finite sum.  On a block $J_l$ with $l<k$, the inductive
formula is simplified in the source proof to
\begin{align}\label{eq:first-inverse}
 \lambda_l(Sh)(\gamma)
  ={}&h(p_{l,l}^{-1}(\gamma))+\lambda_l h(\omega^\alpha+l)\\
    &-h\!\left(p_{l+1,l+1}^{-1}
       \bigl(p_{l,l+1}(p_{l,l}^{-1}(\gamma))\bigr)\right).\nonumber
\end{align}
So each inverse output coordinate uses only three input evaluations.

The second nontrivial map is Proposition 2.15 of~\cite{CHRS}.  Its tuple
length $\ell(\gamma)$ is bounded by a fixed finite integer.  The operator is
given by
\begin{align}\label{eq:second-forward}
 (Tf)(*)&=f(\omega^\alpha),\nonumber\\
 (Tf)(\gamma)
 &=\left(1-\sum_{i=1}^{\ell(\gamma)}\lambda_i\right)f(\omega^\alpha)
  +\sum_{i=1}^{\ell(\gamma)-1}\lambda_i
       f(\rho(\gamma|_i+1))
  +\lambda_{\ell(\gamma)}f(\rho(\gamma)).
\end{align}
Every coordinate is a finite sum.  More decisively, the inverse is explicitly
\begin{align}\label{eq:second-inverse}
 (Sh)(\omega^\alpha)&=h(*),\nonumber\\
 (Sh)(\rho(\gamma))
 &=h(*)+\lambda_1^{-1}(h(\gamma)-h(*))
       &&(\ell(\gamma)=1),\nonumber\\
 (Sh)(\rho(\gamma))
 &=h(*)+\lambda_{\ell(\gamma)}^{-1}
   \bigl(h(\gamma)-h(\gamma|_{\ell(\gamma)-1}+1)\bigr)
       &&(\ell(\gamma)>1).
\end{align}
Hence both directions of this map are coordinate-finite.

Finally, the proof of Theorem 2.4 in~\cite{CHRS} obtains the desired positive
isomorphism by a finite chain consisting of the maps in Propositions 2.12 and
2.15 and homeomorphism-induced isometries.  Lemma~\ref{lem:composition}
therefore gives the asserted coordinate-finiteness of $U$ and $U^{-1}$.
Positivity of $U$ is exactly the positivity verified in the source
construction.
\end{proof}

\section{Full answer}
\begin{theorem}\label{thm:main}
Let $K,L$ be countable compact spaces.  The following are equivalent:
\begin{enumerate}[label=\textup{(\roman*)}]
\item there is a positive linear homeomorphism
      $T:\Cp(K)\to\Cp(L)$;
\item there is a positive Banach-space isomorphism
      $S:\C(K)\to\C(L)$.
\end{enumerate}
If $K,L$ are infinite, these are further equivalent to
\begin{equation}\label{eq:height-conclusion}
 \htop(K)\leq\htop(L)<\Gamma(\htop(K)).
\end{equation}
Thus the answer to Question 4.3 of~\cite{CHRS} is affirmative.
\end{theorem}

\begin{proof}
Assume (i).  Lemma~\ref{lem:closed-graph} makes $T$ and $T^{-1}$ bounded in
the supremum norms.  The algebraic map is already bijective and positive, so
it is the positive Banach isomorphism required in (ii).

Conversely, first suppose $K,L$ are infinite and (ii) holds.  Theorem 2.4 of
\cite{CHRS} says that (ii) is equivalent to
\eqref{eq:height-conclusion}.  Proposition~\ref{prop:source-audit} then
produces a positive Banach isomorphism $U$ for which both $U$ and $U^{-1}$
are coordinate-finite.  Lemma~\ref{lem:cp} says exactly that these two maps
are continuous for the pointwise topologies.  Hence $U:\Cp(K)\to\Cp(L)$ is
the positive linear homeomorphism required in (i).

It remains only to dispose of finite cases.  A finite compact Hausdorff space
is discrete and $\C(K)$ has dimension $|K|$.  If one space is finite and the
other infinite, neither kind of linear isomorphism can exist.  If both are
finite, a Banach-space isomorphism forces $|K|=|L|$; any bijection between
them is a homeomorphism and its composition operator is a positive isometry
and a $C_p$-homeomorphism.  This proves the equivalence in all cases.
\end{proof}

\section{Verification notes and limitations}
\begin{itemize}
\item Lemma~\ref{lem:closed-graph} does not use countability.  Countability is
used only to invoke the source classification and its ordinal construction in
the reverse implication.
\item We do not claim that every positive Banach isomorphism is
$C_p$-continuous.  We prove the existential statement by selecting the
coordinate-finite isomorphism already constructed in the proof of
\cite[Theorem 2.4]{CHRS}.
\item The inverse need not be positive.  Requiring positivity of both a map
and its inverse would be a different order-isomorphism problem.
\item No computational or unproved conditional dependency enters the proof.
The only source-dependent audit is the finite-evaluation form of the explicit
operators and inverses in Propositions 2.12 and 2.15 of~\cite{CHRS}.
\end{itemize}

\section{Bounded novelty search}
On 21 July 2026 we searched the run's lightweight indexes, the local
2000--2026 arXiv source corpus, the current arXiv record for~\cite{CHRS}, and
arXiv web results.  Queries included the exact phrase
\texttt{positive linear homeomorphism C\_p(K)}, the title and arXiv id
\texttt{2601.11463}, all four authors' names, and close phrases involving
positive $C_p$ classification of countable compacta.  The current v1, dated
16 January 2026, still states Question 4.3.  No later explicit answer was
found.  The returned July 2026 work on $C_p$ spaces over separable compact
lines concerns a different classification, and a May 2026 source merely cites
\cite{CHRS} in an operator-space context.

Novelty confidence is moderate rather than definitive.  The implication is
newly identified here, but its short proof relies on observing a property of
the explicit formulas already printed earlier in the same source paper.

\section{Human-review recommendation}
Check the quantifier and orientation in the source proof of Theorem 2.4, then
verify \eqref{eq:first-inverse} and \eqref{eq:second-inverse} directly against
Propositions 2.12 and 2.15.  Those checks establish that every elementary map
and inverse is coordinate-finite.  The remaining arguments are the elementary
$C_p$ dual description, closure under composition, and the closed graph
theorem.

\begin{thebibliography}{9}
\bibitem{CHRS}
M. Cúth, J. Havelka, J. Rondoš, and B. Sarı,
\emph{The classification of $C(K)$ spaces for countable compacta by positive
isomorphisms}, arXiv:2601.11463v1 (2026).
\end{thebibliography}

\end{document}
